\paragraph{What to report.}
The measurements support three practices for any result obtained after an automatic search. (i)~\emph{Disclose the search}: the number of candidates, the number of tuning items, the procedure and its seeds; the synthetic gap is a function of $n$ and $K$ first (Figs.~\ref{fig:gap}, \ref{fig:factors}), and a tuned-vs-untuned or tuned-vs-tuned comparison is not like-for-like unless the budgets are given (Fig.~\ref{fig:claims}). (ii)~\emph{Report on items not used for selection when they can be spared}: select on part of the items and report on a separate set. This is unbiased and needs no modelling of the search. At equal total items it is less accurate than correcting the in-sample score (RMSE $\N{fr_fresh100_200_rmse}$~pp for $100$ tuning plus $100$ fresh items vs.\ $\N{fr_boot_oob_200_rmse}$~pp for the out-of-bag bootstrap at $200$ items, Table~\ref{tab:fresh}), so its price is accuracy, not only width. (iii)~\emph{When items cannot be spared}, either correct the score or use an interval that accounts for the selection. The out-of-bag bootstrap had the smallest error and a small bias in the synthetic reference cell and across the $\N{fr_ncells}$ equal-budget combinations (coverage $\N{fr_oob_covmin}$--$\N{fr_oob_covmax}\%$, mostly conservative), and had bias $\N{real_boot_oob_n200_bias}\pm\N{real_boot_oob_n200_biasse}$~pp, coverage $\N{real_boot_oob_n200_cover}\%$ and, at $200$ items, the smallest RMSE ($\N{frr_boot_oob_200_rmse}$ vs.\ $\N{frr_fresh100_200_rmse}$~pp for $100$/$100$ tuning/fresh items) on the real model ($K=\N{real_K}$); it requires the per-item outcomes of every candidate searched and a search that can be replayed on resampled items, which an adaptive optimiser may not allow. The hybrid winner-aware interval was near-nominal at a width near that of the simultaneous bound in the region where we ran it, but its point estimate was no better than the naive one (RMSE $\N{fr_hybrid_200_rmse}$ vs.\ $\N{fr_naive_200_rmse}$~pp at $200$ items), and the conditional interval was too wide to inform at the sample sizes studied.

\paragraph{Limitations.}
\emph{Synthetic setting.} The generator is a Gaussian copula with a finite configuration pool; its conclusions hold for that generator and the ranges of $n$, $K$, $\tau$, $\rho$ and mean accuracy we scanned (Sec.~\ref{sec:synth}), and the Gaussian-maximum prediction is an approximation that deviates by up to $0.5$~pp. On the real covariance it deviates by at most $\N{gm_maxdiff}$~pp, and a matched generator reproduces the real gap to about ten percent, not exactly (Sec.~\ref{sec:real}). \emph{Real setting.} The pool is $K=\N{real_K}$ configurations with a spread of $\N{real_sd}$~pp, a factorial grid rather than a random sample of prompts, and the searches (best-of-$K$, successive halving, successive rejects, greedy coordinate search) are replayed over that fixed pool rather than run with the LLM in the loop; a real leakage arm, a real two-system comparison and a per-subject analysis were not run, and there is no check at large $K$ and no stronger model. One small model with letter-logit scoring; the truth is the held-out remainder of a finite pool. We did not run DSPy or MIPRO and say nothing about their search dynamics. \emph{Methods.} The conditional and hybrid intervals are our implementations; they pass the checks of Appendix~\ref{app:defs} (closed form, exact-Gaussian coverage, a port of the authors' replication code) but their coverage in the data-generating setting was run only at $\rho=0.5$, mean accuracy $0.3$ and $K\le64$ for the full-evaluation argmax. The equal-budget comparison covers $200$ and $400$ items and best-of-$K$ search; the out-of-bag bootstrap was not run on adaptive search. Cross-fitting estimates the performance of the procedure run on a fraction of the data, not of the final configuration. False-claim and ranking rates depend on the claim threshold and the share of no-better configurations (varied only through $\tau$). Optimistic tuned scores can lead to deploying prompts worse than reported; correct reporting reduces that risk, and we see no direct negative use.
